\section{Introduction}
\label{sec:intro}

MedSAM~\cite{ma2024medsam} adapts the Segment Anything Model~\cite{kirillov2023sam} to medical images and is routinely fine-tuned again for a target modality with LoRA~\cite{hu2022lora} or visual prompt tuning~\cite{jia2022vpt}. A prior comparison on a four-dataset drift ladder showed that the cheapest methods are also the most fragile: LoRA trained on dermoscopy gains in-domain accuracy and then loses 0.18 Dice relative to zero-shot MedSAM on mammography, while decoder-only and full fine-tuning keep their far-OOD accuracy~\cite{prior2026safer}. The same study found that the loss correlates with how far the mask decoder's representations move from the base model, measured by centred kernel alignment (CKA)~\cite{kornblith2019cka}, and that encoder-side movement does not predict it. This paper asks whether that correlation is a mechanism that can be acted on, and what it needs to occur. That fine-tuning distorts pretrained features and hurts out of distribution is established~\cite{kumar2022finetuning,tan2024protection}; the hypothesis tested here is narrower: for a promptable segmenter, far-OOD loss depends on where in the segmentation pathway the representations move, not on how much they move. First, does regularising the decoder toward the base model during adaptation reduce far-OOD collapse, and is the decoder specifically the locus? We add a CKA loss on a 32-image out-of-distribution probe and compare hooks on the decoder, the encoder, and both. Second, does the collapse scale with the adaptation budget? We train three adapters on nested subsets of 50, 250 and 1000 images at a fixed step budget. Third, can far-OOD failures be detected at inference without ground truth? We compare the IoU estimate SAM's decoder already produces with the drift of the decoder's upscaled mask embedding from the base model on the same image and prompt.

Our contributions are:
\begin{enumerate}
  \item CKA-regularised LoRA with an OOD-only probe, three seeds, five jitter levels and a three-seed hook-placement ablation: encoder-only anchoring hurts, decoder anchoring helps, both is best;
  \item a budget study with drift curves: LoRA's far-OOD loss appears at 250 training images with a jump in far-OOD decoder drift, while decoder-only and encoder-only adaptation never collapse or drift, so the failure requires joint, unanchored adaptation;
  \item a label-free per-image failure signal from decoder drift, assessed by threshold sensitivity, risk-coverage and a held-out domain (AUROC 0.96, 0.80 and 0.69 for LoRA over three seeds), and the finding that the decoder's IoU head responds to poor prompts, not to modality shift.
\end{enumerate}
